\section{Topological Spaces}

\begin{definition}[topological space]
Let $X$ be a set and let $\tau\subset \P(X)$ be a family of subsets of $X$.
We will say that $X$ is a \emph{topological space}%
\mymargin{topological space}\index{topological space}, that $\tau$ is a
\emph{topology}%
\mymargin{topology}\index{topology}, and that the elements of $\tau$ are \emph{open}%
\mymargin{open}\index{open}
if:
\begin{enumerate}
\item $\emptyset \in \tau$, $X\in \tau$ (the empty set and the entire space are open);
\item $A \subset \tau \implies \bigcup A \in \tau$ (any union of open sets is open);
\item $A,B\in \tau \implies A \cap B \in \tau$ (finite intersection of open sets is open).
\end{enumerate}
We will also say that the topological space is \emph{separated}%
\mymargin{separated}\index{separated} (or $T_2$ or Hausdorff)
\index{topological space!separated}
\index{topological space!$T_2$}
\index{topological space!Hausdorff}
\index{$T_2$}
\index{separated}
if additionally the following holds:
\begin{enumerate}
\item[4.] given $x,y\in X$ with $x\neq y$, there exist $U,V\in \tau$ with $x\in U$,
$y\in V$ and $U\cap V=\emptyset$.
\end{enumerate}

We will say that a subset $A\subset X$ is \emph{closed}%
\mymargin{closed}\index{closed} if its complement
is open ($X\setminus A\in \tau$).
We will say that $A\subset X$ is a \emph{neighborhood}%
\mymargin{neighborhood}\index{neighborhood} of the point $a\in X$ if there exists $U\in \tau$
such that $x\in U \subset A$.

All the properties given in definition~\ref{def:466342} can thus be
given in a topological space using neighborhoods (or open sets) instead of
balls.
\end{definition}

\begin{theorem}[topology induced by the metric]
Let $X$ be a metric space and let $\tau$ be the set of all open sets
of $X$, that is:
\[
 \tau = \ENCLOSE{A \subset X\colon \forall a \in A \exists \rho>0 \colon B_\rho(a)\subset A}.
\]
Then $\tau$ is a topology and $X$ is a separated topological space with respect
to the topology
$\tau$.
\end{theorem}

\begin{definition}[neighborhood base]
Let $X$ be a set and for each $x\in X$ let $\B_x\subset \P(X)$
be a family of subsets of $X$. We will say that $\B_x$ is a
\emph{neighborhood base}%
\mymargin{neighborhood base}\index{neighborhood base} of $x$ if the following properties hold:
\begin{enumerate}
\item $\B_x$ is non-empty;
\item if $B\in \B_x$ then $x \in B$;
\item if $A,B\in \B_x$ then there exists $C\in \B_x$ such that $C\subset A\cap B$;

\item If a set contains a neighborhood of a point, it is also a neighborhood of that point.
\item The intersection of two neighborhoods of a point is a neighborhood of the point.
\end{enumerate}

If for every $x\in X$ the set $\B_x$ is a neighborhood base of $x$, we can
define
\[
  \tau = \ENCLOSE{A \subset X\colon \forall x\in A \exists B \in \B_x \colon B\subset A}
\]
and it follows that $\tau$ is a topology and thus $X$ is a topological space
with respect to this topology (induced by the neighborhood base).

If additionally the following property holds:
\begin{enumerate}
\item given $x,y\in X$, $x\neq y$, there exist $B\in \B_x$ and $C\in \B_y$ such that
$B\cap C=\emptyset$;
\end{enumerate}
then $X$ will be a separated topological space.
\end{definition}

The following properties are consequences of the previous axioms and are
therefore valid in any topological space.

\begin{enumerate}
\item
The interior of any $A\subset X$ is the largest (i.e., the union of all) open
set contained in $A$. The closure of $A$ is the smallest (i.e., the intersection
of all) closed set containing $A$.
In particular, the interior is always open and the interior of an open set is
the entire set. The closure is a closed set and the closure of a closed set is
the set itself.
\item
The boundary of a set is closed. Interior, boundary, and exterior are three
disjoint sets (respectively open, closed, and open) whose union is the entire
space.
\end{enumerate}

\begin{theorem}
If $X$ is a metric space, the family of sets
\[
  \tau = \ENCLOSE{A \subset X \colon \forall a \in A \exists \rho>0 \colon B_\rho(x) \subset A}
\]
is a topology on $X$ called the \emph{induced topology}
by the metric of $X$.
\end{theorem}

\begin{proof}
The empty set is open because there are no points on which it is necessary to
verify the property that defines open sets. The entire space is also open
because every ball is contained in it. Thus $\emptyset, X$ are open and
consequently their respective complements: $X, \emptyset$ are closed.

That any union of open sets is open is obvious: given a point of the union, such
point is contained in one of the open sets. Thus there is a small ball centered
at the point and contained in the open set. But then it is also contained in the
union.
That the intersection of closed sets is closed is obtained by passing to the
complements: the complement of a closed set is open and the complement of the
intersection is the union of the complements.

The intersection of any closed sets is equal to the complement (with respect to
$X$) of the union of the complements. By definition, the complement of a closed
set is open and thus the union of the complements is open. Therefore, its
complement is closed.

Now consider the intersection $A\cap B$ of two open sets $A$ and $B$. If $x$ is
a point of the intersection, we know there exist two balls $B_r(x)$ and $B_s(x)$
such that $B_r(x)\subset A$ and $B_s(x) \subset B$. The smaller of the two is
contained in $A\cap B$, and this shows that $A\cap B$ is open. Passing to the
complements, it follows that the union of two closed sets is closed.

Given $x,y \in X$, if $x\neq y$, then $r=d(x,y)/2>0$. We then observe that
$B_r(x)$ and $B_r(y)$ are disjoint because if there existed $z\in B_r(x) \cap
B_r(y)$, we would have $d(x,y)\le d(x,z) + d(z,y) < r + r = d(x,y)$, which is
absurd.

We have thus demonstrated the four properties of open sets. Let us move on to
the properties of neighborhoods. By definition, a set $U$ is a neighborhood of
$x$ if there exists an open set $A$ such that $x\in A \subset U$. Necessarily,
therefore, $x\in U$. And if $V\supset U$, then $x\in A \subset V$, and thus $V$
is also a neighborhood of $x$. If $U$ and $V$ are neighborhoods, there must
exist two open sets $A$ and $B$ such that $x\in A \subset U$ and $x \in B
\subset V$. But then $A\cap B$ is open, and $x \in A \cap B \subset U \cap V$,
so $U\cap V$ is also a neighborhood of $x$. For the fourth point, it suffices to
observe that every neighborhood contains, by definition, an open set that
contains the point, and an open set is, again by definition, a neighborhood of
each of its points. For point 5, we already know that two distinct points are
contained in disjoint open sets, and open sets are always neighborhoods of each
of their points.

Let us now demonstrate that the interior of a set $A$ is the union of all open
sets contained in $A$. On one hand, if $x$ is an interior point of $A$, then
there exists $r>0$ such that $B_r(x)\subset A$. Since $B_r(x)$ is open, it
follows that $x$ is in the union of the open sets contained in $A$. Conversely,
if $x$ is in the union of all open sets contained in $A$, there must exist an
open set $U$ such that $x\in U \subset A$. But then there exists $r>0$ such that
$B_r(x) \subset U \subset A$, and thus $x$ is an interior point of $A$. The
analogous property for closed sets is obtained by passing to the complements.

The boundary of a set is, by definition, the set of points that are neither
interior nor exterior. On the other hand, it is clear that a point cannot be
simultaneously interior and exterior. Thus the boundary is the complement of the
union of the interior and the exterior, and is therefore a closed set.
\end{proof}

\begin{theorem}[sequential closure]%
  \label{th:chiuso_sequenziale}%
Let $(X,d)$ be a metric space.
A set $A\subset X$ is closed if and only if
for every sequence $a_n\in A$, if $a_n\to a\in X$, then $a \in A$
(the limit of points of $A$, if it exists, is a point of $A$).
\end{theorem}
%
\begin{proof}
Suppose that $A$ is closed and let
$a_k \in A$ be a sequence converging to a point of $X$: $a_k \to a$.
We must show that $a\in A$.
By definition of convergence,
we know that for every $\eps>0$, there exists $K\in \NN$ such that for every $k>
K$,
we have $d(a_k,a)<\eps$ and thus $a_k \in B_\eps(a)$.
In particular, $A \cap B_\eps(a) \neq \emptyset$.
It follows that $a$ is not exterior to $A$ and therefore, since $A$ is closed,
$a\in A$.

Suppose that $A$ is not closed and verify that in this case
it is possible to find a sequence $a_k$ of points of $A$ that converges $a_k\to
a$
to a point $a\not\in A$.
If $A$ is not closed, it means there is a point $a \in X\setminus A$ that is not
exterior to $A$.
This means that for every $r>0$, the set $B_r(a)\cap A$ is non-empty. For every
$k\in \NN$, I can then
choose $r=1/k$ and thus know that there exists a point $a_k\in B_r(a) \cap A$,
i.e., $a_k \in A$ and
$d(a_k,a) < 1/k$.
Thus $a_k \to a$ with $a_k\in A$ but $a\not \in A$, as we wanted to demonstrate.
\end{proof}

\begin{definition}[continuity]%
  \label{def:continua_metrico}%
\mymark{**}%
A function $f\colon X \to Y$ defined between two metric spaces
is said to be \emph{sequentially continuous}%
\mymargin{sequentially continuous}\index{sequentially!continuous} at the point $x\in X$ if
for every convergent sequence $x_n \to x$ in $X$, 
it follows that $f(x_n)\to f(x)$ in $Y$.
We will say that $f$ is \emph{sequentially continuous} if it is sequentially
continuous
at every point $x$ of its domain $X$.

A function $f\colon X \to Y$ defined between two metric spaces
is said to be \emph{continuous}%
\mymargin{continuous}\index{continuous} at a point $x\in X$ if
\begin{equation}\label{eq:continuita_metrico}
 \forall \eps>0\colon \exists \delta>0\colon \forall y \in X \colon
 d(x,y)< \delta \implies d(f(x), f(y)) < \eps.
\end{equation}
We will say that $f$ is \emph{continuous} if
it is continuous at every point $x$ of its domain.
\end{definition}

Again, we have considered two different notions of continuity that
in general (in topological spaces) might not coincide, but in the case of
metric spaces, they are equivalent, as we demonstrate in the following theorem.

\begin{theorem}[equivalent definitions of continuity]
Let $f\colon X \to Y$ be a function defined between two metric spaces $X$ and
$Y$.
Then $f$ is sequentially continuous at a point $x\in X$ if and only if it is
continuous at the point $x$.

Moreover, $f$ is continuous if and only if
for every $A$ open in $Y$, it follows that $f^{-1}(A)$ is open in $X$
(the preimage of an open set is open).
\end{theorem}
%
\begin{proof}
Suppose that $f$ is sequentially continuous at $x$
and, for contradiction, suppose that $f$ is not continuous at the same point
$x$.
Then by negating~\eqref{eq:continuita_metrico}, we obtain that there exists
$\eps>0$
such that for every $\delta>0$, in particular for every $\delta = \frac{1}{k}$
with $k\in \NN$, there exists a point $y_k$ such that $d(y_k,x)<\frac 1 k$ but
$d(f(y_k),f(x))\ge \eps$. This means that $y_k \to x$ and therefore,
if $f$ is sequentially continuous, it should be $f(y_k)\to f(x)$. But
this is in contradiction with the condition $d(f(y_k),f(x))\ge \eps$.

Conversely, suppose that $f$ is continuous at $x$ and consider any
sequence $x_k \to x$.
To demonstrate that $f(x_k)\to f(x)$, we must verify that for every
$\eps>0$, there exists $K\in \NN$ such that for every $k>K$, we have
$d(f(x_k),f(x)) < \eps$.
But from the continuity of $f$, given $\eps>0$, we know there exists $\delta>0$
such that
if $d(y,x)<\delta$, then $d(f(y),f(x))<\eps$. Since $x_k\to x$,
there certainly exists $K\in \NN$ such that for every $k>K$, we have
$\abs{x_k-x}<\delta$
and thus $\abs{f(x_k)-f(x)}< \eps$, as we wanted to demonstrate.

Now let us show that if $f$ is continuous and $A$ is open in $Y$, then
$f^{-1}(A)$
is open in $X$.
Given any point $x_0\in f^{-1}(A)$, we know that $f(x_0)\in A$, thus, being
$A$ open, there exists $\eps>0$ such that $B_\eps(f(x_0))\subset A$.
By the continuity of $f$, there exists $\delta>0$ such that if $\abs{x-x_0} <
\delta$,
then $\abs{f(x)-f(x_0)}<\eps$ and thus $f(x)\in A$. This means that
$B_\delta(x_0)\subset f^{-1}(A)$.
This is true for every $x_0\in f^{-1}(A)$, so such a set is open.

Conversely, suppose that the preimage of every open set is open and
demonstrate that the function is continuous at every point.
Given a point $x_0\in X$ and an $\eps>0$, consider the open set
$B_\eps(f(x_0))$.
Its preimage is the set $\ENCLOSE{x\colon \abs{f(x)-f(x_0)} < \eps}$ and
by hypothesis, we know it is open. This means that there exists $\delta>0$
such that $B_\delta(x_0)$ is contained in this set, i.e., for every $x\in
B_\delta(x_0)$
that is $\abs{x-x_0}<\delta$,
it follows that $f(x) \in B_\eps(f(x_0))$
that is $\abs{f(x)-f(x_0)}<\eps$.
We have thus verified the definition of
continuity at the point $x_0$.
\end{proof}

\begin{theorem}[continuity of the composite function]%
  \label{th:continuita_funzione_composta_metrico}%
Let $X$, $Y$, and $Z$ be metric spaces, $f\colon X\to Y$ 
and $g\colon Y\to Z$ be continuous functions. 
Then $g\circ f\colon X\to Z$ is continuous.
\end{theorem}
\begin{proof}
  Using sequential continuity, it suffices to observe that 
  if $x_n\to x$, then $f(x_n)\to f(x)$ and therefore $g(f(x_n))\to g(f(x))$.
\end{proof}

\begin{definition}[bounded spaces]%
  \label{def:spazio_limitato}%
\mymark{*}%
Let $X$ be a metric space or a subset of a metric subspace. We will say that $X$
is
\emph{bounded}%
\mymargin{bounded}\index{bounded} if it is contained in a ball, i.e., if
there exists $x_0\in X$ and $R>0$ such that $X\subset B_R(x_0)$.
\end{definition}

\begin{definition}[sequential compactness]%
  \label{def:sequenzialmente_compatto}%
\mymark{**}%
Let $X$ be a metric space or a subset of a
metric space. We will say that $X$ is
\emph{sequentially compact}%
\mymargin{sequentially compact}\index{sequentially!compact} if from every
sequence $x_k \in X$, it is possible to extract a subsequence $x_{k_j}\to x$
converging to a point $x\in X$.
\end{definition}

The more general definition of compactness is given in topological spaces.
In the general context, compactness and sequential compactness are different
concepts,
but in the realm of metric spaces, the two concepts coincide.
For this reason, we can write more briefly \emph{compact}%
\mymargin{compact}\index{compact}
instead of \emph{sequentially compact} when working in metric spaces.

The Bolzano-Weierstrass theorem states that the intervals $[a,b]$ with
$a,b\in \RR$ are compact. 
We can extend this result in $\RR^n$.

\begin{theorem}[compactness in $\RR^n$]%
  \label{th:compattezza_Rn}%
Let $K\subset \RR^n$ be a closed and bounded set. 
Then $K$ is compact.
\end{theorem}
\begin{proof}
Let $\vec x_k\in K$ be any sequence of points. 
If $K$ is bounded, it means there exists a ball $B_r(\vec y)$ 
that contains $K$, thus $\abs{\vec x_k - \vec y} < r$ for every $k$.
Setting $R=r+\abs{\vec y}$, by the reverse triangle inequality, we can assert 
that $\abs{\vec x_k} < R$ for every $k$. 
Denote by $((x_k)_1, \dots, (x_k)_n)$ the coordinates 
of the point $\vec x_k\in \RR^n$. 
Then, since $\abs{(x_k)_j}\le \abs{x_k}$, it follows that $(x_k)_j\in [-R,R]$
for every $k\in\NN$ and every $j=1,\dots, n$.
By the Bolzano-Weierstrass theorem, the sequence $(x_k)_1$ has a convergent
subsequence. But from this subsequence, we can extract a sub-subsequence
for which $(x_k)_2$ is also convergent. 
Obviously, on this sub-subsequence, $(x_k)_1$ continues to be convergent.
Iterating $n$ times, we obtain a sub-sub-\dots-subsequence of $\vec x_k$ 
such that each component $(x_k)_j$ converges to a number $p_j\in \RR$
for $j=1,\dots, n$.
But this means that $\vec x_k \to \vec p \in \RR^n$ with $\vec p = (p_1,\dots,
p_n)$

Since $K$ is closed, it must be $\vec p\in K$, and thus $K$ is compact.
\end{proof}

In a generic metric space, it is not guaranteed that closed and bounded sets 
are compact (a negative example is given by uniform convergence, as
we will see later). 
The reverse implication, however, is always true, as stated in the
following theorem.

\begin{theorem}%
\mymark{**}%
If $A$ is a sequentially compact subset of a metric space $X$,
then $A$ is closed and bounded.
\end{theorem}
%
\begin{proof}
Clearly, $A$ is closed because given a sequence $x_k\in A$ converging to a point
$x\in X$,
we know there exists a subsequence converging to a point of $A$. But
necessarily, every subsequence converges to $x$, so $x\in A$. If $A$ were not
bounded,
fixed $a\in A$, for every $k\in \NN$, there should exist a point $x_k\in A$ such
that $x_k \not\in B_k(a)$,
i.e., $d(x_k,a) > k$. Suppose then that there exists a convergent subsequence
$x_{k_j}\to x \in A$. Then by the reverse triangle inequality, we would have
\[
  d(x, a) \ge d(x_{k_j}, a) - d(x_{k_j},x)
   \ge k_j - d(x_{k_j},x) \to +\infty - 0 = +\infty.
\]
But this is absurd since $d(x,a)\in \RR$.
\end{proof}

\begin{theorem}[Weierstrass: continuous functions map compacts to compacts]
  \label{th:weiestrass_metrico}%
  \index{Weierstrass!theorem of}%
  \index{theorem!of Weierstrass}%
Let $f\colon X \to Y$ be a continuous function between
two metric spaces $X$ and $Y$.
If $K\subset X$ is sequentially compact, then
$f(K)$ is also sequentially compact.
\end{theorem}
%
\begin{proof}
Let $y_k \in f(K)$ be any sequence. Then
there exists $x_k \in K$ such that $f(x_k) = y_k$.
Since $K$ is compact, we can extract a convergent subsequence: $x_{k_j}\to x$.
Since $f$ is continuous, we have
\[
  y_{k_j} = f(x_{k_j}) \to f(x) \in f(K).
\]
\end{proof}

In the case $X=Y=\RR$, we recover the usual Weierstrass theorem, because if
$f\colon [a,b]\to \RR$ is continuous, being $[a,b]$ compact, it follows that
$f([a,b])$ is compact. But the compacts of $\RR$ are closed and bounded, so they
have a maximum and minimum since the supremum and infimum are finite and are
points of adherence of the set.

\begin{definition}[uniform continuity]
    Let $f\colon X \to Y$ be a function defined between two metric spaces $X$
  and $Y$.
  We will say that $f$ is \emph{uniformly continuous}%
\mymargin{uniformly continuous}\index{uniformly continuous} if 
  \[
  \forall \eps>0 \colon \exists \delta>0 \colon 
    d(x,x')< \delta \implies d(f(x),f(x'))<\eps. 
  \]
\end{definition}

\begin{theorem}[Heine-Cantor]
  \mymargin{Heine-Cantor}%
  \index{Heine-Cantor}%
  \index{theorem!of Heine-Cantor}%
  \label{th:HeineCantor}
Let $X$ be a sequentially compact metric space and $Y$ a metric space.
Let $f\colon X\to Y$ be a continuous function. Then $f$ is uniformly continuous.
\end{theorem}
%
\begin{proof}
  The proof is identical to that already given in the corresponding 
  theorem~\ref{th:heine_cantor} for functions of a real variable.
\end{proof}